% Part: computability
% Chapter: recursive-functions
% Section: composition

\documentclass[../../../include/open-logic-section]{subfiles}

\begin{document}

\olfileid[hi]{cmp}{rec}{com}
\olsection{संयोजन}

यदि $f$ और $g$ प्राकृतिक संख्याओं के दो एक-स्थानी फलन हैं, तो हम उनका
संयोजन कर सकते हैं: $h(x) = f(g(x))$। तब नया फलन~$h(x)$, फलनों $f$
और~$g$ से \emph{संयोजन} द्वारा परिभाषित होता है। हम इसे एक से अधिक
तर्क वाले फलनों के लिए सामान्यीकृत करना चाहते हैं।

ऐसा करने का एक ढंग यह है: मान लें कि $f$ कोई $k$-स्थानी फलन है और
$g_0$, \dots, $g_{k-1}$ ऐसे $k$ फलन हैं जो सभी $n$-स्थानी हैं। तब हम
एक नया $n$-स्थानी फलन $h$ इस प्रकार परिभाषित कर सकते हैं:
\[
h(x_0, \dots, x_{n-1}) =
f(g_0(x_0, \dots, x_{n-1}), \dots, g_{k-1}(x_0, \dots, x_{n-1}))
\]
यदि $f$ और सभी $g_i$ संगणनीय हैं, तो $h$ भी संगणनीय है। $h(x_0,
\dots, x_{n-1})$ संगणित करने के लिए पहले मान $y_i = g_i(x_0, \dots,
x_{n-1})$ को प्रत्येक $i = 0$, \dots,~$k-1$ के लिए संगणित करें।
फिर इन मानों को $f$ में देकर
$h(x_0, \dots, x_{k-1}) = f(y_0, \dots, y_{k-1})$ संगणित करें।

कुछ विद्यमान फलनों का उपयोग करके नया फलन संगणित करने पर जो होता है, यह
उसका अत्यधिक प्रतिबंधित अभिलक्षणन लग सकता है। एक बात तो यह है कि कभी-कभी
हम किसी फलन के सभी तर्कों का उपयोग नहीं करते, जैसे हमने
$g(x, y, z) = \Succ(z)$ को~$\Add$ की आदिम पुनरावर्ती परिभाषा में उपयोग
के लिए परिभाषित किया था। मान लें कि हमें निम्न फलनों का उपयोग करने की अनुमति है:
\[
\Proj{n}{i}(x_0, \dots, x_{n-1}) = x_i
\]
फलनों~$\Proj{k}{i}$ को \emph{प्रक्षेपण} फलन कहते हैं:
$\Proj{n}{i}$ एक $n$-स्थानी फलन है। तब $g$ को इस प्रकार परिभाषित किया जा सकता है:
\[
g(x, y, z) = \Succ(\Proj{3}{2}(x, y, z)).
\]
यहाँ $f$ की भूमिका $1$-स्थानी फलन $\Succ$ निभाता है, इसलिए $k=1$ है।
और हमारे पास एक $3$-स्थानी फलन $\Proj{3}{2}$ है जो $g_0$ की भूमिका
निभाता है। परिणाम एक ऐसा $3$-स्थानी फलन है जो तीसरे तर्क का उत्तराधिकारी
लौटाता है।

प्रक्षेपण फलन तर्कों का क्रम बदलकर या उनकी पहचान करके भी नए फलन परिभाषित
करने देते हैं। मसलन, फलन $h(x)
= \Add(x, x)$ को इस प्रकार परिभाषित किया
जा सकता है:
\[
h(x_0) = \Add(\Proj{1}{0}(x_0),\Proj{1}{0}(x_0)).
\]
यहाँ $k=2$, $n=1$ है, $f(y_0,y_1)$ की भूमिका $\Add$ निभाता है, और
$g_0(x_0)$ तथा $g_1(x_0)$ दोनों की भूमिका~$\Proj{1}{0}(x_0)$ निभाता
है, जो एक-स्थानी प्रक्षेपण फलन (अर्थात् तत्समक फलन) है।

यदि $f(y_0, y_1)$ हमारे पास पहले से कोई फलन है, तो हम फलन
$h(x_0, x_1) = f(x_1, x_0)$ को इस प्रकार परिभाषित कर सकते हैं:
\[
h(x_0, x_1) = f(\Proj{2}{1}(x_0, x_1),\Proj{2}{0}(x_0, x_1)).
\]
यहाँ $k=2$, $n = 2$ है, और $g_0$ तथा $g_1$ की भूमिकाएँ क्रमशः
$\Proj{2}{1}$ और~$\Proj{2}{0}$ निभाते हैं।

आपको यह चिंता भी हो सकती है कि $g_0$, \dots,~$g_{k-1}$ सभी की
स्थान-संख्या समान~$n$ होना आवश्यक है। (याद रखें कि किसी फलन की
\emph{स्थान-संख्या} उसके तर्कों की संख्या है; किसी $n$-स्थानी फलन की
स्थान-संख्या~$n$ है।) किंतु प्रक्षेपण फलनों को जोड़ने से वांछित लचीलापन
मिलता है। मसलन, मान लें कि $f$ और~$g$, $3$-स्थानी फलन हैं तथा $h$ वह
$2$-स्थानी फलन है जिसे इस प्रकार परिभाषित किया गया है:
\[
h(x,y) = f(x,g(x,x,y),y).
\]
प्रक्षेपण फलनों के साथ~$h$ की परिभाषा को इस प्रकार फिर लिखा जा सकता है:
\[
h(x,y) = f(\Proj{2}{0}(x,y), g(\Proj{2}{0}(x,y), \Proj{2}{0}(x,y),
\Proj{2}{1}(x,y)), \Proj{2}{1}(x,y)).
\]
तब $h$, $f$ का $\Proj{2}{0}$, $l$ और $\Proj{2}{1}$ के साथ संयोजन है, जहाँ
\[
l(x,y) = g(\Proj{2}{0}(x,y),\Proj{2}{0}(x,y),\Proj{2}{1}(x,y)),
\]
अर्थात् $l$, $g$ का $\Proj{2}{0}$, $\Proj{2}{0}$ और~$\Proj{2}{1}$ के
साथ संयोजन है।

\end{document}
